\documentclass[11pt]{article}

\usepackage[utf8]{inputenc}
\usepackage[T1]{fontenc}
\usepackage{graphicx}
\usepackage{caption}
\usepackage{subcaption}
\usepackage{hyperref}
\usepackage[margin=1in]{geometry}
\usepackage{xcolor}
\usepackage{float}

% Figures are stored alongside this file under output/run_N/.
\graphicspath{{./}{output/run_9/}{output/run_4/}{output/run_1/}}

% Helper for visible placeholders the successor must resolve.
\newcommand{\fillin}[1]{\textcolor{red}{\textbf{[FILL: #1]}}}
\newcommand{\figplaceholder}[1]{%
  \begin{center}
  \fbox{\parbox{0.7\linewidth}{\centering\vspace{1.5cm}%
  \textbf{[FIGURE PLACEHOLDER]}\\[0.5em]#1\vspace{1.5cm}}}
  \end{center}}

\title{Waste Detection Workaround: A CPU Pipeline Combining\\
SAM 2 and Grounding DINO}
\author{\fillin{author} \\ \fillin{affiliation}}
\date{May 29, 2026}

\begin{document}

\maketitle

\begin{abstract}
This report documents an abandoned prototype for waste segmentation in
photographs. The pipeline combines Meta's SAM 2 segmentation model with
Grounding DINO (referred to in project notes as ``dyno'') for open-vocabulary
detection, and runs entirely on CPU under an explicit memory budget. A
free-text prompt is converted into bounding boxes by Grounding DINO, and each
box is segmented into a precise mask by SAM 2. The approach functioned as a
working baseline. The repository is preserved as both a reusable CPU pipeline
and a documented record of prompt and threshold tuning. We describe the
methods, present qualitative outputs from the recorded runs, assess quality
where evidence exists, and outline how the scaffolding can be carried forward by
a successor. No quantitative ground truth is available, so quality assessment is
qualitative and clearly marked as such.
\end{abstract}

\section{Introduction}
The objective of this work was to optimise object detection and classification
of waste images by combining two pretrained models in a two-stage,
detect-then-segment design. The setting imposes two practical constraints
visible in the codebase: execution is restricted to CPU, and memory must be
budgeted explicitly to avoid exhausting available RAM. The prototype was
intended as a workaround under these constraints rather than a production
system. This document records what was built and what it produced.

\section{Methods}

\subsection{Pipeline overview}
The single entry point, \texttt{sam2\_cpu\_batch.py}, reads every image in the
\texttt{input/} folder, processes it according to the selected mode, and writes
colored mask overlays to a fresh \texttt{output/run\_N/} folder together with a
\texttt{parameters.txt} snapshot of the settings used. The data flow is
detect (or grid-sample), segment, upscale, suppress overlaps, and overlay.

\figplaceholder{Pipeline diagram: input image $\rightarrow$ resize $\rightarrow$
Grounding DINO (boxes) $\rightarrow$ SAM 2 (masks) $\rightarrow$ upscale
$\rightarrow$ mask-IoU NMS $\rightarrow$ overlay $\rightarrow$ run folder.
No rendered diagram file exists yet.}

\subsection{Models}
Two detection modes are active in the current script:
\begin{itemize}
  \item \textbf{\texttt{gdino} (default).} Grounding DINO
  (\texttt{IDEA-Research/grounding-dino-tiny} or \texttt{grounding-dino-base})
  converts a text prompt into bounding boxes, and \texttt{SAM2ImagePredictor}
  segments each box.
  \item \textbf{\texttt{grid}.} \texttt{SAM2AutomaticMaskGenerator} segments via
  a dense point grid with no text prompt.
\end{itemize}
SAM 2.1 is available in three checkpoint sizes present in the repository:
\texttt{hiera\_tiny}, \texttt{hiera\_base\_plus}, and \texttt{hiera\_large}. A
third detection mode, \texttt{florence2}, appears only in the legacy logs of
\texttt{run\_2} and \texttt{run\_3} and is no longer part of the script.

\subsection{Pre- and post-processing}
\begin{itemize}
  \item \textbf{Resize} (\texttt{resize\_for\_inference}): images are scaled so
  the longest side does not exceed \texttt{max\_image\_size}, bounding peak
  memory.
  \item \textbf{Mask upscale} (\texttt{upscale\_masks}): masks computed on the
  resized image are nearest-neighbour upscaled back to the original resolution.
  \item \textbf{Mask-IoU NMS} (\texttt{suppress\_overlapping\_masks}): when two
  masks overlap above \texttt{mask\_nms\_iou\_threshold}, the larger one is
  discarded so the tighter mask is kept.
  \item \textbf{Overlay} (\texttt{overlay\_masks}): masks are blended onto the
  image with deterministic, seed-controlled random colors.
\end{itemize}
Execution is bounded by \texttt{reserved\_cpu\_cores} (PyTorch threads) and a
\texttt{psutil}-based RAM pre-check that warns when the estimated peak exceeds
the usable share of free memory.

\subsection{Configuration history}
The recorded runs trace the evolution of the approach:
\begin{itemize}
  \item \textbf{run\_1}: \texttt{grid} baseline with \texttt{hiera\_tiny},
  \texttt{points\_per\_side=40}, \texttt{pred\_iou\_thresh=0.6}.
  \item \textbf{run\_2, run\_3}: legacy \texttt{florence2} mode with
  \texttt{microsoft/Florence-2-base-ft} (no overlays retained; mode since
  removed).
  \item \textbf{run\_4--run\_9}: \texttt{gdino} mode, progressively tuning the
  detection prompt, detection thresholds, model size, and
  \texttt{max\_image\_size}.
\end{itemize}
The final recorded configuration (run\_9) used the \texttt{grounding-dino-base}
detector, the \texttt{hiera\_base\_plus} segmenter, the multi-class prompt
\texttt{construction waste. general waste. plastic tubes. wood. plastic.},
\texttt{max\_image\_size=2048}, \texttt{threshold=0.25}, and
\texttt{text\_threshold=0.1}.

\section{Results}
All results below are qualitative. The repository contains no labelled ground
truth, so no precision, recall, or IoU-against-truth figures can be computed
from the current artifacts.

\subsection{Final tuned run (gdino, run\_9)}
\begin{figure}[H]
  \centering
  \begin{subfigure}[b]{0.48\linewidth}
    \includegraphics[width=\linewidth]{output/run_9/image00010_overlay.png}
    \caption{\texttt{image00010}}
  \end{subfigure}
  \hfill
  \begin{subfigure}[b]{0.48\linewidth}
    \includegraphics[width=\linewidth]{output/run_9/image00020_overlay.png}
    \caption{\texttt{image00020}}
  \end{subfigure}
  \caption{Mask overlays from the final tuned \texttt{gdino} run (run\_9):
  \texttt{grounding-dino-base} + SAM 2.1 \texttt{hiera\_base\_plus}, multi-class
  prompt, \texttt{max\_image\_size=2048}, \texttt{text\_threshold=0.1}.}
  \label{fig:run9}
\end{figure}

\subsection{Early versus tuned comparison}
\begin{figure}[H]
  \centering
  \begin{subfigure}[b]{0.48\linewidth}
    \includegraphics[width=\linewidth]{output/run_4/image00020_overlay.png}
    \caption{Early \texttt{gdino} (run\_4): single-class prompt
    \texttt{construction waste.}, \texttt{hiera\_tiny}.}
  \end{subfigure}
  \hfill
  \begin{subfigure}[b]{0.48\linewidth}
    \includegraphics[width=\linewidth]{output/run_9/image00020_overlay.png}
    \caption{Tuned \texttt{gdino} (run\_9): multi-class prompt,
    \texttt{hiera\_base\_plus}, larger input.}
  \end{subfigure}
  \caption{Qualitative effect of tuning on \texttt{image00020}: prompt
  enrichment, a larger detector and segmenter, and a higher input resolution.}
  \label{fig:compare}
\end{figure}

\subsection{Quality assessment}
\figplaceholder{Quantitative quality chart (e.g. detection count or
mask-coverage per run). No metric data or ground-truth annotations exist on
disk; this figure cannot be produced without a labelled evaluation set.}

Per-run configuration deltas observed in the recorded \texttt{parameters.txt}
files:
\begin{itemize}
  \item \textbf{Prompt} grew from \texttt{construction waste.} (run\_4) to a
  five-class prompt (run\_9), broadening the categories detected.
  \item \textbf{Detector} moved from \texttt{grounding-dino-tiny} to
  \texttt{grounding-dino-base}.
  \item \textbf{Segmenter} moved from \texttt{hiera\_tiny} to
  \texttt{hiera\_base\_plus}.
  \item \textbf{Input resolution} (\texttt{max\_image\_size}) increased from
  1024 to 2048.
  \item \textbf{\texttt{text\_threshold}} was lowered from 0.25 to 0.1,
  admitting weaker text-to-region matches.
\end{itemize}

\section{Discussion}
\paragraph{Strengths.}
The CLI scaffolding is model-agnostic: resource budgeting, run-folder
management, mask-IoU NMS, and deterministic overlay rendering are independent of
the chosen detector and segmenter. Execution is resource-aware, with an explicit
CPU-thread reservation and a RAM pre-check. Every run is self-documenting via
its \texttt{parameters.txt} snapshot.

\paragraph{Weaknesses.}
Quality is judged only qualitatively; there is no evaluation harness or labelled
data. Results are sensitive to the detection prompt and thresholds. CPU
execution is slow for the larger checkpoints, particularly
\texttt{hiera\_large}. The removed \texttt{florence2} branch is undocumented and
survives only in legacy logs.

\paragraph{Project status.}
Development on this prototype has stopped and the repository is being handed off.
The SAM 2 + Grounding DINO pipeline is retained as a baseline and as reusable
scaffolding rather than as a finished production solution.

\section{Conclusions}
The repository is a viable CPU baseline for prompt-driven waste segmentation and
a documented log of the exploration carried out so far. Its surrounding
pipeline---resource management, per-run logging, NMS, and overlays---remains
useful independently of the specific models, and is the most directly reusable
asset for a successor.

\section{Future Work}
\begin{itemize}
  \item Introduce a labelled evaluation set to enable quantitative quality
  metrics and principled threshold sweeps.
  \item Add a non-interactive run mode (config file or CLI flags) for automated
  batches.
  \item Add a pinned dependency file and an initial git commit to complete the
  handoff.
\end{itemize}

\section*{References and Acknowledgements}
\begin{itemize}
  \item \fillin{SAM 2 citation}
  \item \fillin{Grounding DINO citation}
\end{itemize}

\end{document}
